% Independent replacement for Section 3; the manuscript entry point is unchanged.
% The matching appendix is defined below as \LIFTNewMethodProofs, to be invoked
% in the appendix when this draft is integrated. No other source is modified.
% Integration issues: the existing abstract, introduction, experiments and
% appendix still describe the previous lens/barycentric method.
% Implementation issue: restore_with_budget in minimal_validity_repair.py can
% move initially invalid supports beyond their first-valid points. The method
% below specifies the requested minimal restoration followed by rejection of
% over-budget candidates; the additional implementation fallback is not minimal.
% The current driver expands a training attack bank before lambda selection.
% This is construction-time calibration, not independent final evaluation.

\section{Distributionally Robust Counterfactual Refinement}
\label{sec:new-method}

We propose \textsc{LIFT-CF}, a post-hoc refinement of the population produced by a base CF generator, optimizing worst-case future invalidity under model updates.
Unrestricted robust optimization can move this population arbitrarily far from both the generator's output and the favorable target population.
We constrain the total transport through these two anchors, allowing a bounded detour beyond their shortest Wasserstein distance.
For computation, we optimize continuous support coordinates against a finite bank of feasible model updates, then restore current validity only where needed.
We first quantify the price of the geometric restriction and then describe the numerical realization and deployment safeguard.

\subsection{Problem Setting and Deployment Robustness}
\label{sec:new-objective}

Let $f_0:\mathcal X\to[0,1]$ be the current classifier and $\tau\in(0,1)$ the decision threshold.
An input $x$ is favorable if $f_0(x)\ge\tau$.
For unfavorable instances $X=\{x_i\}_{i=1}^n$, a base generator $G$ produces $z_i=G(x_i)$, defining $\mathbb P_{\mathrm{CF}}=n^{-1}\sum_{i=1}^n\delta_{z_i}$.
Let $\mathbb P_1$ denote the empirical distribution of favorable training samples.
These anchors describe the generator's placement of CFs and the observed favorable-class structure, respectively.

Write $\widehat{\mathbb P}_{\mathrm{tr}}$ for the empirical training distribution and $W_{\mathrm{data}}$ for the data-side Wasserstein metric.
For radius $r\ge0$, define $\mathcal U_r=\{\mathbb P:W_{\mathrm{data}}(\mathbb P,\widehat{\mathbb P}_{\mathrm{tr}})\le r\}$ and $f_{\mathbb P}=\operatorname{Update}(f_0,\mathbb P)$, where the update procedure is fixed.
With invalidity loss $\ell_f(z)=[\tau-f(z)]_+$, the exact deployment robust risk is
\begin{equation}
\Phi_r(\mathbb Q)
:=\sup_{\mathbb P\in\mathcal U_r}
\mathbb E_{z\sim\mathbb Q}[\ell_{f_{\mathbb P}}(z)].
\label{eq:new-robust-risk}
\end{equation}
The supremum selects one updated model for the entire CF population.
It is this population risk that we minimize, while leaving $G$ unchanged.

\subsection{Bounded Transport-Detour Decision Region}
\label{sec:new-geometry}

A counterfactual is a displacement: an input is moved to a favorable state.
Refining a population of CFs therefore naturally induces a transport problem.
Wasserstein distance lifts the underlying feature-space modification cost to distributions and remains meaningful when their empirical supports do not overlap.
Let $d_{\mathrm{CF}}$ denote this ground metric and $W_{\mathrm{CF}}$ its $2$-Wasserstein metric on $\mathcal P_2(\mathcal X)$, the distributions with finite second moment.
We assume this Wasserstein space is geodesic, as for Euclidean features on a closed convex domain.

Restricting refinement to a single shortest geodesic can exclude robust distributions that require a detour.
Instead, let $D:=W_{\mathrm{CF}}(\mathbb P_{\mathrm{CF}},\mathbb P_1)$ and, for a geometry budget $B_{\mathrm{geo}}\ge D$, define
\begin{equation}
\mathcal G_{B_{\mathrm{geo}}}
:=\left\{\mathbb Q\in\mathcal P_2(\mathcal X):
W_{\mathrm{CF}}(\mathbb P_{\mathrm{CF}},\mathbb Q)
+W_{\mathrm{CF}}(\mathbb Q,\mathbb P_1)\le B_{\mathrm{geo}}\right\}.
\label{eq:new-region}
\end{equation}
Our central refinement problem is
\begin{equation}
\boxed{\inf_{\mathbb Q\in\mathcal G_{B_{\mathrm{geo}}}}\Phi_r(\mathbb Q)
=\inf_{\mathbb Q\in\mathcal G_{B_{\mathrm{geo}}}}
\sup_{\mathbb P\in\mathcal U_r}
\mathbb E_{z\sim\mathbb Q}[\ell_{f_{\mathbb P}}(z)].}
\label{eq:new-wro}
\end{equation}
Here $r$ controls uncertainty in future training data, whereas $B_{\mathrm{geo}}$ controls geometric freedom on the CF side.
For $D>0$, the normalized budget is $\beta_{\mathrm{geo}}:=B_{\mathrm{geo}}/D\ge1$.

This restriction has a direct geometric interpretation.
The geodesic excess $e(\mathbb Q):=W_{\mathrm{CF}}(\mathbb P_{\mathrm{CF}},\mathbb Q)+W_{\mathrm{CF}}(\mathbb Q,\mathbb P_1)-D$ is the additional transport detour relative to the shortest connection between the anchors.
Thus, $\mathcal G_{B_{\mathrm{geo}}}$ allows excess at most $B_{\mathrm{geo}}-D$ in exchange for robustness.
This geometric departure is distinct from paired recourse cost: it does not measure the increase in distance from each factual input to its assigned CF.
We report that downstream cost separately.

\begin{proposition}[Zero-detour special case]
\label{prop:new-zero-detour}
In the stated geodesic space, $\mathcal G_D$ is exactly the union of the shortest-geodesic segments between $\mathbb P_{\mathrm{CF}}$ and $\mathbb P_1$.
If $D=0$, it consists only of their common anchor.
\end{proposition}
The triangle inequality forces equality when $B_{\mathrm{geo}}=D$, recovering the shortest-geodesic construction as the zero-detour case; it does not select a unique geodesic.

\subsection{Robustness under Geometric Restriction}
\label{sec:new-theory}

We next quantify the price of this geometric restriction.
Assume the invalidity losses are uniformly Lipschitz with respect to $d_{\mathrm{CF}}$, with $L_r:=\sup_{\mathbb P\in\mathcal U_r}\operatorname{Lip}(\ell_{f_{\mathbb P}})<\infty$.
Then $|\Phi_r(\mathbb Q)-\Phi_r(\mathbb Q')|\le L_rW_{\mathrm{CF}}(\mathbb Q,\mathbb Q')$.
For any budget $B\ge D$, write $V_r(B):=\inf_{\mathbb Q\in\mathcal G_B}\Phi_r(\mathbb Q)$.

\begin{corollary}[Nested decision regions]
\label{cor:new-nested}
For $D\le B_1\le B_2$, $\mathcal G_{B_1}\subseteq\mathcal G_{B_2}$ and $V_r(B_2)\le V_r(B_1)$.
\end{corollary}

\begin{theorem}[Price of geometric restriction]
\label{thm:new-restriction}
Suppose the uniform Lipschitz condition holds and an unrestricted minimizer $\mathbb Q_r^\star\in\arg\min_{\mathbb Q\in\mathcal P_2(\mathcal X)}\Phi_r(\mathbb Q)$ exists.
Define its representation distance by $\rho_r(B):=\inf_{\mathbb Q\in\mathcal G_B}W_{\mathrm{CF}}(\mathbb Q,\mathbb Q_r^\star)$.
For every $B\ge D$,
\begin{equation}
0\le V_r(B)-\Phi_r(\mathbb Q_r^\star)\le L_r\rho_r(B).
\label{eq:new-restriction-bound}
\end{equation}
\end{theorem}
Since $\rho_r(B)$ is nonincreasing in $B$, additional geometric freedom can only reduce this representation restriction.
The bound concerns the distributional decision region, not global convergence of its numerical solver.
Proofs are given in Appendix~\ref{app:new-method-proofs}.

\subsection{Finite-Bank Continuous Realization}
\label{sec:new-realization}

For computation, we represent $\mathbb Q=\sum_{j=1}^m b_j\delta_{w_j}$ with fixed masses $b_j\ge0$, $\sum_j b_j=1$, and continuous support coordinates $w_j\in\mathcal X$.
We replace the inner supremum by a finite bank $\mathcal A_r\subseteq\{f_{\mathbb P}:\mathbb P\in\mathcal U_r\}$ of feasible full-update models:
\begin{equation}
\widehat\Phi_r(\mathbb Q)
:=\max_{f\in\mathcal A_r}\mathbb E_{z\sim\mathbb Q}[\tau-f(z)]_+
=\max_{f\in\mathcal A_r}\sum_{j=1}^m b_j[\tau-f(w_j)]_+.
\label{eq:new-bank-risk}
\end{equation}
The aggregation is the maximum over models of the population mean hinge, not the mean of each CF's individual worst-model loss.
Adversarial searches propose shifted training distributions; retained models are obtained by the fixed full update and checked for data-side feasibility.
Consequently, $\widehat\Phi_r(\mathbb Q)\le\Phi_r(\mathbb Q)$ is an adversarial lower approximation to the inner supremum.

We optimize the support coordinates under the single constraint in Eq.~\ref{eq:new-region}, without a current-validity penalty or a restriction to graph nodes.
Numerical candidates are assessed using the population bank hinge and recomputed Wasserstein distances.
Finite particle representation, finite-bank coverage, and local numerical optimization are separate approximations to Eq.~\ref{eq:new-wro}.
After freezing the construction, fresh adversarial searches expand the bank for evaluation; they do not retroactively tune the final test result.

\subsection{Minimal Validity Restoration and Deployment}
\label{sec:new-deployment}

The raw robust supports need not be favorable under $f_0$.
For each invalid support $w_j^{\mathrm{raw}}$, let $y_j\in\operatorname{supp}(\mathbb P_1)$ be its matched favorable target anchor and set $w_j(\alpha)=(1-\alpha)w_j^{\mathrm{raw}}+\alpha y_j$.
Minimal validity restoration chooses $\alpha_j^\star:=\min\{\alpha\in[0,1]:f_0(w_j(\alpha))\ge\tau\}$ and replaces the support by $w_j(\alpha_j^\star)$.
For continuous $f_0$ and segments contained in $\mathcal X$, this minimum exists because the anchor is favorable.
Already-valid supports remain unchanged; invalid supports move only as far as needed along their matched segments, without assuming monotonic scores along a segment.
We recompute both transport distances after restoration and retain a repaired candidate only if their sum remains at most $B_{\mathrm{geo}}$.
This is a numerical deployment safeguard, separate from Theorem~\ref{thm:new-restriction}; it does not imply preservation of robust risk.

\paragraph{Operating-point realization.}
The implementation retains a path through the learned control distribution $\mathbb Q$.
Gluing optimal couplings from $\mathbb P_{\mathrm{CF}}$ to $\mathbb Q$ and from $\mathbb Q$ to $\mathbb P_1$ gives a piecewise displacement path $\mathbb Q(\lambda)$, parameterized by normalized arc length $\lambda\in[0,1]$.
This path realizes deployment operating points after learning $\mathbb Q$; it is not the theoretical decision set.
On a predefined grid, we select the current-valid population with smallest mean-hinge risk against a shared training attack bank, breaking exact ties toward smaller $\lambda$.
Training-bank expansion used for this selection remains part of construction.
The final operating point is frozen before held-out evaluation, and its validity and geometry are checked explicitly; no held-out future outcomes choose $\lambda$.

% Supporting appendix, kept in this file to respect the requested edit scope.
% Invoke \LIFTNewMethodProofs after \appendix when integrating the draft.
\newcommand{\LIFTNewMethodProofs}{%
\section{Proofs for Geometry-Constrained Refinement}
\label{app:new-method-proofs}

\begin{proof}[Proof of Proposition~\ref{prop:new-zero-detour}]
For $\mathbb Q\in\mathcal G_D$, the triangle inequality gives
$D\le W_{\mathrm{CF}}(\mathbb P_{\mathrm{CF}},\mathbb Q)+W_{\mathrm{CF}}(\mathbb Q,\mathbb P_1)\le D$.
Thus equality holds.
In a geodesic space, concatenate shortest segments from the first anchor to $\mathbb Q$ and from $\mathbb Q$ to the second anchor.
Their total length is $D$, so the concatenation is shortest: any shorter subsegment would yield an anchor-to-anchor path shorter than $D$.
Parameterizing by arc length places $\mathbb Q$ on a constant-speed shortest geodesic.
Conversely, every point on such a geodesic has distances to its endpoints summing to $D$ and belongs to $\mathcal G_D$.
If $D=0$, both distances must vanish, giving the common anchor.
\end{proof}

\begin{proof}[Proof of Corollary~\ref{cor:new-nested}]
Every distribution satisfying the transport-sum constraint with budget $B_1$ satisfies it with budget $B_2\ge B_1$.
Taking an infimum of the same objective over this larger set cannot increase its value.
\end{proof}

\begin{proof}[Proof of Theorem~\ref{thm:new-restriction}]
For every $\mathbb P\in\mathcal U_r$ and every coupling of $\mathbb Q,\mathbb Q'$, the uniform Lipschitz condition bounds the difference of expected losses by $L_r$ times the expected ground distance.
Infimizing over couplings gives a bound by the corresponding $1$-Wasserstein distance, which is at most $W_{\mathrm{CF}}$ by Cauchy--Schwarz.
Taking suprema over $\mathbb P$ in both directions yields
$|\Phi_r(\mathbb Q)-\Phi_r(\mathbb Q')|\le L_rW_{\mathrm{CF}}(\mathbb Q,\mathbb Q')$.
Unrestricted optimality gives $V_r(B)\ge\Phi_r(\mathbb Q_r^\star)$.
For each $\mathbb Q\in\mathcal G_B$, we also have
$V_r(B)\le\Phi_r(\mathbb Q)\le\Phi_r(\mathbb Q_r^\star)+L_rW_{\mathrm{CF}}(\mathbb Q,\mathbb Q_r^\star)$.
Infimizing the right-hand side over $\mathcal G_B$ proves the upper bound, without requiring attainment of the representation distance.
Nestedness also shows that $\rho_r(B)$ is nonincreasing.
\end{proof}
}
